\documentclass[10pt,a4paper]{amsart}
\usepackage{amsmath,amssymb,amsthm,graphicx}
% Retained minimal preamble for the Stacks Project chapter varieties.tex.
% Verbatim excerpt of the chapter preamble: Xy-pic setup (native preamble lines 14-18), the native theorem
% declarations with the mathematical macros (native preamble lines 154-217), and the project own no-op
% definitions of the reference/slogan/history display environments (native wrapper.tex lines 1-3) that the
% counted statement uses for its slogan. Copied mechanically; nothing added, removed or redefined.
% For commutative diagrams we use Xy-pic
\usepackage[all]{xy}

% We use 2cell for 2-commutative diagrams.
\xyoption{2cell}
% Theorem environments.
%
\theoremstyle{plain}
\newtheorem{theorem}[subsection]{Theorem}
\newtheorem{proposition}[subsection]{Proposition}
\newtheorem{lemma}[subsection]{Lemma}

\theoremstyle{definition}
\newtheorem{definition}[subsection]{Definition}
\newtheorem{example}[subsection]{Example}
\newtheorem{exercise}[subsection]{Exercise}
\newtheorem{situation}[subsection]{Situation}

\theoremstyle{remark}
\newtheorem{remark}[subsection]{Remark}
\newtheorem{remarks}[subsection]{Remarks}

\numberwithin{equation}{subsection}

% Macros
%
\def\lim{\mathop{\mathrm{lim}}\nolimits}
\def\colim{\mathop{\mathrm{colim}}\nolimits}
\def\Spec{\mathop{\mathrm{Spec}}}
\def\Hom{\mathop{\mathrm{Hom}}\nolimits}
\def\Ext{\mathop{\mathrm{Ext}}\nolimits}
\def\SheafHom{\mathop{\mathcal{H}\!\mathit{om}}\nolimits}
\def\SheafExt{\mathop{\mathcal{E}\!\mathit{xt}}\nolimits}
\def\Sch{\mathit{Sch}}
\def\Mor{\mathop{\mathrm{Mor}}\nolimits}
\def\Ob{\mathop{\mathrm{Ob}}\nolimits}
\def\Sh{\mathop{\mathit{Sh}}\nolimits}
\def\NL{\mathop{N\!L}\nolimits}
\def\CH{\mathop{\mathrm{CH}}\nolimits}
\def\proetale{{pro\text{-}\acute{e}tale}}
\def\etale{{\acute{e}tale}}
\def\QCoh{\mathit{QCoh}}
\def\Ker{\mathop{\mathrm{Ker}}}
\def\Im{\mathop{\mathrm{Im}}}
\def\Coker{\mathop{\mathrm{Coker}}}
\def\Coim{\mathop{\mathrm{Coim}}}

% Boxtimes
%
\DeclareMathSymbol{\boxtimes}{\mathbin}{AMSa}{"02}

%
% Macros for moduli stacks/spaces
%
\def\QCohstack{\mathcal{QC}\!\mathit{oh}}
\def\Cohstack{\mathcal{C}\!\mathit{oh}}
\def\Spacesstack{\mathcal{S}\!\mathit{paces}}
\def\Quotfunctor{\mathrm{Quot}}
\def\Hilbfunctor{\mathrm{Hilb}}
\def\Curvesstack{\mathcal{C}\!\mathit{urves}}
\def\Polarizedstack{\mathcal{P}\!\mathit{olarized}}
\def\Complexesstack{\mathcal{C}\!\mathit{omplexes}}
% \Pic is the operator that assigns to X its picard group, usage \Pic(X)
% \Picardstack_{X/B} denotes the Picard stack of X over B
% \Picardfunctor_{X/B} denotes the Picard functor of X over B
\def\Pic{\mathop{\mathrm{Pic}}\nolimits}
\def\Picardstack{\mathcal{P}\!\mathit{ic}}
\def\Picardfunctor{\mathrm{Pic}}
\def\Deformationcategory{\mathcal{D}\!\mathit{ef}}
\newenvironment{reference}{}{}
\newenvironment{slogan}{}{}
\newenvironment{history}{}{}
\usepackage{cleveref}
\crefname{originalreference}{original reference}{original references}
\Crefname{originalreference}{Original reference}{Original references}
\makeatletter
\newlabel{algebra-lemma-Noetherian-irreducible-components}{{00FR}{source}{}{sourcepointer-1}{}}
\newlabel{algebra-lemma-Noetherian-irreducible-components@cref}{{[originalreference][0][]00FR}{[1][0][]source}{}{}{}}
\makeatother
\crefname{originalreference}{original reference}{original references}
\Crefname{originalreference}{Original reference}{Original references}
\makeatletter
\newlabel{algebra-lemma-Zariski-topology}{{00E0}{source}{}{sourcepointer-2}{}}
\newlabel{algebra-lemma-Zariski-topology@cref}{{[originalreference][0][]00E0}{[1][0][]source}{}{}{}}
\makeatother
\crefname{originalreference}{original reference}{original references}
\Crefname{originalreference}{Original reference}{Original references}
\makeatletter
\newlabel{algebra-lemma-spec-localization}{{00E3}{source}{}{sourcepointer-3}{}}
\newlabel{algebra-lemma-spec-localization@cref}{{[originalreference][0][]00E3}{[1][0][]source}{}{}{}}
\makeatother
\crefname{originalreference}{original reference}{original references}
\Crefname{originalreference}{Original reference}{Original references}
\makeatletter
\newlabel{algebra-proposition-dimension-zero-ring}{{00KJ}{source}{}{sourcepointer-4}{}}
\newlabel{algebra-proposition-dimension-zero-ring@cref}{{[originalreference][0][]00KJ}{[1][0][]source}{}{}{}}
\makeatother
\crefname{originalreference}{original reference}{original references}
\Crefname{originalreference}{Original reference}{Original references}
\makeatletter
\newlabel{algebra-lemma-height-1}{{00KK}{source}{}{sourcepointer-5}{}}
\newlabel{algebra-lemma-height-1@cref}{{[originalreference][0][]00KK}{[1][0][]source}{}{}{}}
\makeatother
\crefname{originalreference}{original reference}{original references}
\Crefname{originalreference}{Original reference}{Original references}
\makeatletter
\newlabel{algebra-lemma-bound-depth}{{00LK}{source}{}{sourcepointer-6}{}}
\newlabel{algebra-lemma-bound-depth@cref}{{[originalreference][0][]00LK}{[1][0][]source}{}{}{}}
\makeatother
\crefname{originalreference}{original reference}{original references}
\Crefname{originalreference}{Original reference}{Original references}
\makeatletter
\newlabel{algebra-lemma-criterion-reduced}{{031R}{source}{}{sourcepointer-7}{}}
\newlabel{algebra-lemma-criterion-reduced@cref}{{[originalreference][0][]031R}{[1][0][]source}{}{}{}}
\makeatother
\crefname{originalreference}{original reference}{original references}
\Crefname{originalreference}{Original reference}{Original references}
\makeatletter
\newlabel{algebra-proposition-minimal-primes-associated-primes}{{02CE}{source}{}{sourcepointer-8}{}}
\newlabel{algebra-proposition-minimal-primes-associated-primes@cref}{{[originalreference][0][]02CE}{[1][0][]source}{}{}{}}
\makeatother
\crefname{originalreference}{original reference}{original references}
\Crefname{originalreference}{Original reference}{Original references}
\makeatletter
\newlabel{algebra-lemma-ass-zero-divisors}{{00LD}{source}{}{sourcepointer-9}{}}
\newlabel{algebra-lemma-ass-zero-divisors@cref}{{[originalreference][0][]00LD}{[1][0][]source}{}{}{}}
\makeatother
\crefname{originalreference}{original reference}{original references}
\Crefname{originalreference}{Original reference}{Original references}
\makeatletter
\newlabel{algebra-lemma-total-ring-fractions-no-embedded-points}{{02LX}{source}{}{sourcepointer-10}{}}
\newlabel{algebra-lemma-total-ring-fractions-no-embedded-points@cref}{{[originalreference][0][]02LX}{[1][0][]source}{}{}{}}
\makeatother
\crefname{originalreference}{original reference}{original references}
\Crefname{originalreference}{Original reference}{Original references}
\makeatletter
\newlabel{algebra-lemma-minimal-prime-reduced-ring}{{00EU}{source}{}{sourcepointer-11}{}}
\newlabel{algebra-lemma-minimal-prime-reduced-ring@cref}{{[originalreference][0][]00EU}{[1][0][]source}{}{}{}}
\makeatother
\title{Equivalent characterizations of an element of the maximal ideal that cuts out the maximal ideal up to radical}
\author{The Stacks Project Authors}
\date{stacks-project:varieties.tex:local-snapshot-20261011; GNU FDL 1.2 or later}
\begin{document}
\maketitle
\noindent\textit{Editorial note.} All delivered mathematical text is the author's own, copied verbatim from the pinned Stacks Project chapter file varieties.tex; no statement or proof step was rewritten, reordered or removed.\par\medskip
\noindent\textit{Editorial note.} The counted span cites eleven cross-chapter labels (including propositions), each recorded in assembly.external\_references with its canonical tag as a navigation pointer.\par\medskip
\noindent\textit{Editorial note.} The proof cites no bibliography item inside the counted span.\par\medskip
\noindent\textit{Allowed standard prerequisites (editorial).}\begin{itemize}
\item Zariski topology criterion for the radical of a principal ideal; spectrum of a localization. Stacks Project, Algebra, Lemma algebra-lemma-Zariski-topology (tag 00IY) and Lemma algebra-lemma-spec-localization (tag 00E4). Give the equivalence of (1) and (2) and the description of Spec(A\_f).
\item Zero-dimensional rings as products of localizations; associated primes; total ring of fractions. Stacks Project, Algebra, Proposition algebra-proposition-dimension-zero-ring (tag 00JA); Proposition algebra-proposition-minimal-primes-associated-primes (tag 00LD); Lemma algebra-lemma-total-ring-fractions-no-embedded-points (tag 02LX). Supply the identification of A\_f with a product of the localizations and the depth-one and total-ring-of-fractions equivalences.
\end{itemize}\medskip
\ifdefined\hypertarget\hypertarget{sourcepointer-1}{}\fi
\noindent\textit{Source reference.} Original reference 00FR: sources/stacks-chow/tags.txt, source lines 581--581. Cross-chapter standard prerequisite cited by the author inside the counted span; it is used as a pointer and its own proof is not part of this item.\par\medskip
\ifdefined\hypertarget\hypertarget{sourcepointer-2}{}\fi
\noindent\textit{Source reference.} Original reference 00E0: sources/stacks-chow/tags.txt, source lines 517--517. Cross-chapter standard prerequisite cited by the author inside the counted span; it is used as a pointer and its own proof is not part of this item.\par\medskip
\ifdefined\hypertarget\hypertarget{sourcepointer-3}{}\fi
\noindent\textit{Source reference.} Original reference 00E3: sources/stacks-chow/tags.txt, source lines 520--520. Cross-chapter standard prerequisite cited by the author inside the counted span; it is used as a pointer and its own proof is not part of this item.\par\medskip
\ifdefined\hypertarget\hypertarget{sourcepointer-4}{}\fi
\noindent\textit{Source reference.} Original reference 00KJ: sources/stacks-chow/tags.txt, source lines 753--753. Cross-chapter standard prerequisite cited by the author inside the counted span; it is used as a pointer and its own proof is not part of this item.\par\medskip
\ifdefined\hypertarget\hypertarget{sourcepointer-5}{}\fi
\noindent\textit{Source reference.} Original reference 00KK: sources/stacks-chow/tags.txt, source lines 754--754. Cross-chapter standard prerequisite cited by the author inside the counted span; it is used as a pointer and its own proof is not part of this item.\par\medskip
\ifdefined\hypertarget\hypertarget{sourcepointer-6}{}\fi
\noindent\textit{Source reference.} Original reference 00LK: sources/stacks-chow/tags.txt, source lines 790--790. Cross-chapter standard prerequisite cited by the author inside the counted span; it is used as a pointer and its own proof is not part of this item.\par\medskip
\ifdefined\hypertarget\hypertarget{sourcepointer-7}{}\fi
\noindent\textit{Source reference.} Original reference 031R: sources/stacks-chow/tags.txt, source lines 3994--3994. Cross-chapter standard prerequisite cited by the author inside the counted span; it is used as a pointer and its own proof is not part of this item.\par\medskip
\ifdefined\hypertarget\hypertarget{sourcepointer-8}{}\fi
\noindent\textit{Source reference.} Original reference 02CE: sources/stacks-chow/tags.txt, source lines 3081--3081. Cross-chapter standard prerequisite cited by the author inside the counted span; it is used as a pointer and its own proof is not part of this item.\par\medskip
\ifdefined\hypertarget\hypertarget{sourcepointer-9}{}\fi
\noindent\textit{Source reference.} Original reference 00LD: sources/stacks-chow/tags.txt, source lines 783--783. Cross-chapter standard prerequisite cited by the author inside the counted span; it is used as a pointer and its own proof is not part of this item.\par\medskip
\ifdefined\hypertarget\hypertarget{sourcepointer-10}{}\fi
\noindent\textit{Source reference.} Original reference 02LX: sources/stacks-chow/tags.txt, source lines 3424--3424. Cross-chapter standard prerequisite cited by the author inside the counted span; it is used as a pointer and its own proof is not part of this item.\par\medskip
\ifdefined\hypertarget\hypertarget{sourcepointer-11}{}\fi
\noindent\textit{Source reference.} Original reference 00EU: sources/stacks-chow/tags.txt, source lines 548--548. Cross-chapter standard prerequisite cited by the author inside the counted span; it is used as a pointer and its own proof is not part of this item.\par\medskip
\begin{lemma}
\label{lemma-pre-pre-delta-invariant}
Let $(A, \mathfrak m)$ be a Noetherian $1$-dimensional local ring.
Let $f \in \mathfrak m$. The following are equivalent
\begin{enumerate}
\item $\mathfrak m = \sqrt{(f)}$,
\item $f$ is not contained in any minimal prime of $A$, and
\item $A_f = \prod_{\mathfrak p\text{ minimal}} A_\mathfrak p$ as $A$-algebras.
\end{enumerate}
Such an $f \in \mathfrak m$ exists. If $\text{depth}(A) = 1$ (for example
$A$ is reduced), then (1) -- (3) are also equivalent to
\begin{enumerate}
\item[(4)] $f$ is a nonzerodivisor,
\item[(5)] $A_f$ is the total ring of fractions of $A$.
\end{enumerate}
If $A$ is reduced, then (1) -- (5) are also equivalent to
\begin{enumerate}
\item[(6)] $A_f$ is the product of the residue fields at the minimal
primes of $A$.
\end{enumerate}
\end{lemma}

\begin{proof}
The spectrum of $A$ has finitely many primes
$\mathfrak p_1, \ldots, \mathfrak p_n$ besides $\mathfrak m$
and these are all minimal, see
Algebra, Lemma \ref{algebra-lemma-Noetherian-irreducible-components}.
Then the equivalence of (1) and (2) follows from
Algebra, Lemma \ref{algebra-lemma-Zariski-topology}.
Clearly, (3) implies (2). Conversely, if (2) is true,
then the spectrum of $A_f$ is the subset
$\{\mathfrak p_1, \ldots, \mathfrak p_n\}$ of $\Spec(A)$
with induced topology, see
Algebra, Lemma \ref{algebra-lemma-spec-localization}.
This is a finite discrete topological space.
Hence $A_f = \prod_{\mathfrak p\text{ minimal}} A_\mathfrak p$ by
Algebra, Proposition \ref{algebra-proposition-dimension-zero-ring}.
The existence of an $f$ is asserted in
Algebra, Lemma \ref{algebra-lemma-height-1}.

\medskip\noindent
Assume $A$ has depth $1$. (This is the maximum by
Algebra, Lemma \ref{algebra-lemma-bound-depth} and holds if $A$ is reduced by
Algebra, Lemma \ref{algebra-lemma-criterion-reduced}.)
Then $\mathfrak m$ is not an associated prime of $A$.
Every minimal prime of $A$ is an associated prime
(Algebra, Proposition
\ref{algebra-proposition-minimal-primes-associated-primes}).
Hence the set of nonzerodivisors of $A$ is exactly the set of elements
not contained in any of the minimal primes by
Algebra, Lemma \ref{algebra-lemma-ass-zero-divisors}.
Thus (4) is equivalent to (2).
Part (5) is equivalent to (3) by
Algebra, Lemma \ref{algebra-lemma-total-ring-fractions-no-embedded-points}.

\medskip\noindent
Then $A_\mathfrak p$ is a field for
$\mathfrak p \subset A$ minimal, see
Algebra, Lemma \ref{algebra-lemma-minimal-prime-reduced-ring}.
Hence (3) is equivalent to (6).
\end{proof}

\end{document}
